% Generated draft: requires mathematical and source review.
\documentclass[11pt]{article}
\usepackage[margin=25mm]{geometry}
\usepackage{fontspec}
\setmainfont{Latin Modern Roman}
\usepackage{amsmath,amssymb,mathtools,unicode-math}
\setmathfont{Latin Modern Math}
\usepackage{enumitem,needspace,xurl,hyperref}
\setlength{\parindent}{0pt}
\setlength{\parskip}{5pt}
\setlength{\emergencystretch}{3em}
\newcommand{\N}{\mathbb{N}}
\newcommand{\Z}{\mathbb{Z}}
\newcommand{\Q}{\mathbb{Q}}
\newcommand{\R}{\mathbb{R}}
\newcommand{\C}{\mathbb{C}}
\newcommand{\F}{\mathbb{F}}
\newcommand{\A}{\mathbb{A}}
\newcommand{\PP}{\mathbb{P}}
\newcommand{\E}{\mathbb{E}}
\newcommand{\eps}{\varepsilon}
\newcommand{\dd}{\,\mathrm{d}}
\begin{document}
\section*{UnsolvedMath ID: 30004883}
\textbf{Source problem code:} catalogue entry
\subsection{Definitions and mathematical statement}
Let \(Q\) be the quiver in the cited source, let \(Q_0\) denote its vertex set, and let \(\Pi_Q\) denote its preprojective algebra. A dimension vector is a tuple \(d=(d_i)_{i\in Q_0}\in\mathbb N^{Q_0}\), recording the dimensions at the vertices. The base field, hypotheses on \(Q\), and convention for \(\mathbb N\) are not specified in the supplied excerpt and must be recovered from the source.

Write \(A_{\Pi_Q}\) for the preprojective cohomological Hall algebra mentioned in the record, and \(\mathfrak g_{\Pi_Q,d}\) for the degree-\(d\) generator object occurring in its quoted Theorem 3. Let \(I_d\) denote the object written \(\operatorname{IC}(X_d)\) in the clean statement. Here \(\operatorname{IC}\) denotes an intersection complex; the space \(X_d\), its coefficients, normalization, and the passage from this complex to a space of generators are not defined in the supplied record.

\textbf{Conjectural completeness target.} For a dimension vector \(d\) such that a simple \(\Pi_Q\)-module of dimension vector \(d\) exists, the supplied account of Theorem 3 describes a canonical inclusion of the intersection-complex contribution into the degree-\(d\) primitive-generator object. Subject to the source's precise identification of these objects, write this inclusion as
\[
\iota_d:I_d\hookrightarrow\mathfrak g_{\Pi_Q,d}.
\]
The intended completeness assertion is that this canonical inclusion exhausts the primitive degree-\(d\) generators:
\[
\operatorname{im}(\iota_d)=\mathfrak g_{\Pi_Q,d},
\qquad\text{equivalently, }\iota_d\text{ is an isomorphism}.
\]
This is a conditional reconstruction of the target, not a fully specified conjecture: the supplied text does not define the generator object or the meaning of primitive, and does not establish that \(\mathfrak g_{\Pi_Q,d}\) is literally a primitive subspace of a specified coalgebra. The original excerpt also contains a corrupted canonical direct-sum decomposition with an intersection-complex term involving \(\mathbb Q^{\mathrm{vir}}\); its exact coefficients, shifts, and complementary summand cannot be recovered. The simple-module existence hypothesis is explicitly attached to Theorem 3; whether the conjecture is intended to include other dimension vectors remains unresolved.
\subsection{Short English statement}
The conjecture asks whether the canonical contribution denoted \(\operatorname{IC}(X_d)\) accounts for every primitive generator of degree \(d\) in the preprojective cohomological Hall algebra. The supplied theorem discussion assumes that a simple module of dimension vector \(d\) exists. Definitions of \(X_d\), the generator object, and the relevant conventions are missing, so the displayed formulation requires source consultation and human mathematical review.
\subsection{Sources}
\begin{itemize}
\item Supplied UnsolvedMath local registry record, permanent ID \(30004883\), title \emph{Primitive Generators of Preprojective Cohomological Hall Algebras}, exact source problem code \texttt{catalogue entry}. The record gives proposed year 2021 and names no proposer.
\item Supplied report citation: \emph{Enveloping Algebras and Geometric Representation Theory} (2022), Oberwolfach Reports, DOI 10.4171/owr/2021/52, pp. 2827--2891. Source URL and literature-source link labeled ``OWR report'': \url{https://doi.org/10.4171/owr/2021/52}.
\item Supplied literature-source link labeled ``Davison'': \url{https://arxiv.org/abs/2007.03289}. The record supplies no further bibliographic details or precise location for the relevant claim.
\item Literature assessment dated \textbf{2026-08-22}: the supplied curation labels the problem ``partially solved'' and reports that Theorem 3 constructs a canonical intersection-complex contribution, while completeness remains conjectural. This is supplied dated curation, not an independently verified assessment of the literature or present status.
\item The registry labels the statement \texttt{corrected\_verified} and reports that the official Oberwolfach report was checked directly. That verification claim is supplied provenance; it has not been independently verified for this draft.
\end{itemize}
\end{document}
